\documentclass[9pt,aspectratio=169]{beamer}
\usetheme{StAndrews}

\coursename{FI5699 Dissertation Module}

\title{\LARGE \textnormal{Week 5: \\
Key Concepts and Terms}}
\hypersetup{pdftitle={FI5699 Week 5 -- Key Concepts and Terms}}

\newcolumntype{R}{>{\raggedright\arraybackslash}p{8cm}}

\begin{document}

% ============================================================================
% TITLE PAGE
% ============================================================================
\begin{frame}[plain]
\titlepage
\end{frame}


\begin{frame}{How to Use This Study Aid}

This deck collects the \alert{key concepts and terms} from Week 5: Panel Regression Recap.

\medskip

\textbf{For each term you should be able to:}
\begin{itemize}
    \item Define it in your own words
    \item Give an example or illustration
    \item Explain why it matters for your dissertation workflow
\end{itemize}

\medskip

Use these slides as prompts --- prepare your own definitions, answers and materials.

\medskip

\alert{Note:} This is a study aid only. It does not attempt to be comprehensive and should not be relied upon as sufficient preparation for assessment purposes.

\end{frame}


% ============================================================================
% SECTION 1: PANEL DATA STRUCTURE
% ============================================================================
\section{Panel Data Structure}

\begin{frame}{What Is Panel Data?}

\begin{tabular}{@{}lR@{}}
\toprule
\textbf{Term} & \textbf{What to know} \\
\midrule
\alert{Panel data} & What are its two dimensions? How does it differ from pure cross-sectional or pure time-series data? \\[0.8em]
\alert{Cross-section} & What does the cross-sectional dimension represent (e.g.\ firms, countries)? \\[0.8em]
\alert{Time series} & What does the time dimension represent (e.g.\ years, months)? \\[0.8em]
\alert{Unit--time pair} & What is a ``firm--year'' observation? Why is each row identified by two things? \\
\bottomrule
\end{tabular}

\end{frame}


\begin{frame}{Why Panel Data Matters}

\begin{tabular}{@{}lR@{}}
\toprule
\textbf{Term} & \textbf{What to know} \\
\midrule
\alert{Unobserved heterogeneity} & What are examples of things you cannot directly measure (e.g.\ firm culture, management quality)? \\[0.8em]
\alert{Within-unit variation} & What does it mean to compare a firm ``to itself'' over time? \\[0.8em]
\alert{Between-unit variation} & What does it mean to compare different firms to each other? Why can this be misleading? \\
\bottomrule
\end{tabular}

\end{frame}


% ============================================================================
% SECTION 2: THE POOLED OLS PROBLEM
% ============================================================================
\section{The Pooled OLS Problem}

\begin{frame}{Pooled OLS and Omitted Variable Bias}

\begin{tabular}{@{}lR@{}}
\toprule
\textbf{Term} & \textbf{What to know} \\
\midrule
\alert{Pooled OLS} & What does it mean to run OLS on panel data without fixed effects? What assumption does it make? \\[0.8em]
\alert{Omitted variable bias (OVB)} & What happens to $\hat{\beta}$ when an unobserved variable correlates with both $X$ and $Y$? \\[0.8em]
\alert{Confounding variable} & What is a confounder? How does it create a spurious relationship between $X$ and $Y$? \\[0.8em]
\alert{Endogeneity} & What does it mean for $X$ to be correlated with the error term $\varepsilon$? \\
\bottomrule
\end{tabular}

\end{frame}


% ============================================================================
% SECTION 3: FIXED EFFECTS
% ============================================================================
\section{Fixed Effects}

\begin{frame}{The Fixed Effects Idea}

\begin{tabular}{@{}lR@{}}
\toprule
\textbf{Term} & \textbf{What to know} \\
\midrule
\alert{Fixed effect ($\alpha_i$)} & What is a unit-specific intercept? What does it absorb? \\[0.8em]
\alert{Demeaning} & What does it mean to subtract each unit's time average? Why is this equivalent to including $\alpha_i$? \\[0.8em]
\alert{Time-invariant} & What kinds of variables are ``absorbed'' by firm fixed effects (e.g.\ founding year, geography)? \\[0.8em]
\alert{Singleton observation} & What happens to firms with only one observation? Why are they dropped? \\
\bottomrule
\end{tabular}

\end{frame}


\begin{frame}{Types of Fixed Effects}

\begin{tabular}{@{}lR@{}}
\toprule
\textbf{Term} & \textbf{What to know} \\
\midrule
\alert{Firm FE ($\alpha_i$)} & What time-invariant firm traits does it absorb? \\[0.8em]
\alert{Year FE ($\delta_t$)} & What common time shocks does it absorb (e.g.\ financial crises, regulatory changes)? \\[0.8em]
\alert{Two-way FE} & What does including both firm and year FE simultaneously control for? \\[0.8em]
\alert{Industry $\times$ Year FE} & What does an interacted fixed effect absorb that separate FEs cannot? \\
\bottomrule
\end{tabular}

\end{frame}


\begin{frame}{Limitations of Fixed Effects}

\begin{tabular}{@{}lR@{}}
\toprule
\textbf{Term} & \textbf{What to know} \\
\midrule
\alert{No dynamics assumption} & Why does FE require that past outcomes do not cause current treatment? (Imai \& Kim 2019) \\[0.8em]
\alert{Within variation only} & Why does $\hat{\beta}$ come only from within-unit changes, not cross-sectional differences? \\[0.8em]
\alert{Plausible counterfactuals} & Why should you use within-unit SD (not overall SD) for effect-size claims? (Mummolo \& Peterson 2018) \\[0.8em]
\alert{Residualizing $X$} & What does it mean to ``residualize'' $X$ with respect to your fixed effects, and why does it reveal the variation driving $\hat{\beta}$? \\
\bottomrule
\end{tabular}

\end{frame}


% ============================================================================
% SECTION 4: CLUSTERED STANDARD ERRORS
% ============================================================================
\section{Clustered Standard Errors}

\begin{frame}{Why OLS Standard Errors Fail in Panels}

\begin{tabular}{@{}lR@{}}
\toprule
\textbf{Term} & \textbf{What to know} \\
\midrule
\alert{Independence assumption} & What does OLS assume about the correlation between $\varepsilon_{it}$ and $\varepsilon_{js}$? Why does this fail in panels? \\[0.8em]
\alert{Firm effect (serial correlation)} & Why are residuals for the same firm correlated across years? \\[0.8em]
\alert{Time effect (cross-sectional correlation)} & Why are residuals for different firms correlated within the same year? \\[0.8em]
\alert{Over-rejection} & What happens to hypothesis tests when SEs are too small? (Petersen 2009: 15\% reject at 1\%) \\
\bottomrule
\end{tabular}

\end{frame}


\begin{frame}{Clustering Solutions}

\begin{tabular}{@{}lR@{}}
\toprule
\textbf{Term} & \textbf{What to know} \\
\midrule
\alert{Clustered standard errors} & What assumption do they relax? How do they allow correlation within clusters? \\[0.8em]
\alert{Cluster by firm} & When should you cluster by firm? What dependence structure does it handle? \\[0.8em]
\alert{Cluster by time} & When should you cluster by time? When might you prefer Fama--MacBeth? \\[0.8em]
\alert{Two-way clustering} & What does $\hat{V}_{\text{firm}} + \hat{V}_{\text{time}} - \hat{V}_{\text{White}}$ achieve? When do you need it? \\[0.8em]
\alert{Minimum cluster count} & Why do you need at least 30--50 clusters for reliable inference? \\
\bottomrule
\end{tabular}

\end{frame}


\begin{frame}{When to Cluster}

\begin{tabular}{@{}lR@{}}
\toprule
\textbf{Term} & \textbf{What to know} \\
\midrule
\alert{Petersen's diagnostic} & How does comparing SE methods (IID, robust, firm-clustered, time-clustered) reveal the dependence structure? \\[0.8em]
\alert{Treatment assignment level} & Why should you cluster at (or above) the level at which treatment varies? \\[0.8em]
\alert{Over-clustering} & What do Abadie et al.\ (2023) show about unnecessary clustering destroying statistical power? \\[0.8em]
\alert{Robust (HC1) SEs} & When might heteroskedasticity-robust SEs suffice instead of clustering? \\
\bottomrule
\end{tabular}

\end{frame}


% ============================================================================
% SECTION 5: RUNNING REGRESSIONS WITH PYFIXEST
% ============================================================================
\section{Running Regressions with \texttt{pyfixest}}

\begin{frame}{The Formula Syntax}

\begin{tabular}{@{}lR@{}}
\toprule
\textbf{Term} & \textbf{What to know} \\
\midrule
\alert{\texttt{pf.feols()}} & What are its three key arguments (formula, vcov, data)? \\[0.8em]
\alert{Formula: \texttt{y \textasciitilde{} x1 + x2}} & What does the left-hand side represent? What does \texttt{+} mean on the right? \\[0.8em]
\alert{The \texttt{|} separator} & What goes after the pipe? How does it differ from Stata's \texttt{absorb()}? \\[0.8em]
\alert{\texttt{vcov} argument} & What are the options: \texttt{"iid"}, \texttt{"HC1"}, \texttt{\{"CRV1": "firm"\}}, \texttt{\{"CRV1": "firm + year"\}}? \\
\bottomrule
\end{tabular}

\end{frame}


\begin{frame}{Extracting and Comparing Results}

\begin{tabular}{@{}lR@{}}
\toprule
\textbf{Term} & \textbf{What to know} \\
\midrule
\alert{\texttt{.coef()}, \texttt{.se()}, \texttt{.pvalue()}} & What does each method return from a fitted model? \\[0.8em]
\alert{\texttt{.tidy()}} & What does it return and why is it useful for further analysis? \\[0.8em]
\alert{\texttt{.confint()}} & What is a confidence interval and how do you extract one? \\[0.8em]
\alert{\texttt{.vcov()} after estimation} & Why is it useful to re-compute SEs without re-running the regression? How does this implement Petersen's diagnostic? \\[0.8em]
\alert{\texttt{csw0()}} & What does cumulative stepwise estimation do? Why run models that progressively add controls? \\
\bottomrule
\end{tabular}

\end{frame}


% ============================================================================
% SECTION 6: PRESENTING RESULTS
% ============================================================================
\section{Presenting Results}

\begin{frame}{Regression Tables}

\begin{tabular}{@{}lR@{}}
\toprule
\textbf{Term} & \textbf{What to know} \\
\midrule
\alert{\texttt{pf.etable()}} & How does it display multiple models side by side? What arguments does it take? \\[0.8em]
\alert{\texttt{labels}} & Why should you replace variable names with human-readable labels? \\[0.8em]
\alert{\texttt{type="tex"}} & How do you export a table as \LaTeX{} for your dissertation? \\[0.8em]
\alert{\texttt{coef\_fmt}} & How do you control what appears below the coefficient (SEs vs $t$-stats)? \\
\bottomrule
\end{tabular}

\end{frame}


\begin{frame}{Coefficient Plots and Hypothesis Testing}

\begin{tabular}{@{}lR@{}}
\toprule
\textbf{Term} & \textbf{What to know} \\
\midrule
\alert{\texttt{pf.coefplot()}} & What does a coefficient plot show? When is it more effective than a table? \\[0.8em]
\alert{Confidence interval} & What does the interval around each coefficient represent? \\[0.8em]
\alert{\texttt{hypotheses()} (marginaleffects)} & How do you test whether two coefficients are equal? \\[0.8em]
\alert{Delta method} & What does it compute and why is it needed for non-linear functions of coefficients (e.g.\ ratios)? \\
\bottomrule
\end{tabular}

\end{frame}


% ============================================================================
% SECTION 7: KEY PAPERS
% ============================================================================
\section{Key Papers}

\begin{frame}{Four Papers You Should Know}

\begin{tabular}{@{}lR@{}}
\toprule
\textbf{Paper} & \textbf{Key insight} \\
\midrule
\alert{Petersen (2009)} & Compare SE methods to diagnose the dependence structure of your data; firm effect $\to$ cluster by firm. \\[0.8em]
\alert{Imai \& Kim (2019)} & FE assumes no dynamic feedback between treatment and outcome; violations make $\hat{\beta}$ inconsistent. \\[0.8em]
\alert{Mummolo \& Peterson (2018)} & Use within-unit SD for counterfactuals; many published effect sizes are overstated. \\[0.8em]
\alert{Abadie et al.\ (2023)} & Cluster at the level of treatment assignment; unnecessary clustering wastes power. \\
\bottomrule
\end{tabular}

\end{frame}


% ============================================================================
% SECTION 8: SUMMARY
% ============================================================================
\section{Summary Checklist}

\begin{frame}{Week 5 --- Key Terms Checklist}

\textbf{Can you define and give an example for each?}

\medskip

\begin{columns}[T]
\begin{column}{0.33\textwidth}
\begin{itemize}
    \item Panel data
    \item Cross-section / time series
    \item Unit--time pair
    \item Pooled OLS
    \item Omitted variable bias
    \item Endogeneity
\end{itemize}
\end{column}
\begin{column}{0.33\textwidth}
\begin{itemize}
    \item Fixed effect ($\alpha_i$)
    \item Demeaning
    \item Firm / year / two-way FE
    \item Industry $\times$ year FE
    \item Singleton observation
    \item Within-unit variation
\end{itemize}
\end{column}
\begin{column}{0.33\textwidth}
\begin{itemize}
    \item Clustered standard errors
    \item Two-way clustering
    \item Petersen's diagnostic
    \item \texttt{pf.feols()} / the \texttt{|} separator
    \item \texttt{pf.etable()} / \texttt{pf.coefplot()}
    \item Delta method
\end{itemize}
\end{column}
\end{columns}

\end{frame}


\begin{frame}{Week 5 --- Key Questions}

\textbf{Can you answer these in your own words?}

\begin{enumerate}
    \item What is panel data and how does it differ from cross-sectional data?
    \item What is omitted variable bias and why does pooled OLS suffer from it?
    \item What does a firm fixed effect absorb, and what can it \textit{not} absorb?
    \item Why did the patent stock coefficient drop from 0.847 to 0.486 when firm FE were added?
    \item What is the trade-off when adding more fixed effects to a model?
    \item Why are OLS standard errors biased in panel data?
    \item When should you cluster by firm vs.\ by time vs.\ two-way?
    \item In \texttt{pyfixest}, what does the \texttt{|} in the formula \texttt{y \textasciitilde{} x | firm + year} do?
\end{enumerate}

\end{frame}


\end{document}
